\begin{center}
{\large\textbf{Diffraction due to surface tension waves on water}}\footnote{\textit{Shirish Pathare (HBCSE, Mumbai) and K G M Nair (CMI, Chennai) were the principal authors of this problem. The contributions of the Academic Committee, Academic Development Group and the International Board are gratefully acknowledged.}}
\end{center}

\textbf{Introduction}

Formation and propagation of waves on a liquid surface are important and
well-studied phenomena. For such waves, the restoring force on the oscillating
liquid is partly due to gravity and partly due to surface tension. For
wavelengths much smaller than a critical wavelength, $\lambda_\mathrm{c}$, the
effect of gravity is negligible and only surface tension effects need be
considered ($\lambda_c = 2\pi\sqrt{\frac{\sigma}{\rho g}}$, where $\sigma$ is
the surface tension, $\rho$ is the density of the liquid and $g$ is the
acceleration due to gravity).

In this part, you will study surface tension waves on the surface of a liquid,
which have wavelengths smaller than $\lambda_\mathrm{c}$. Surface tension is a
property of liquids due to which the liquid surface behaves like a stretched
membrane. When the liquid surface is disturbed, the disturbance propagates as a
wave just as on a membrane. An electrically-driven vibrator is used to produce
waves on the water surface. When a laser beam is incident at a glancing angle
on these surface waves, they act as a reflection grating, producing a
well-defined diffraction pattern.

Surface tension waves are damped (their amplitude gradually decreases) as they
propagate. This damping is due to the viscosity of the liquid, a property where
adjacent layers of a liquid oppose relative motion between them.

\textbf{Objective}

To use diffraction from surface tension waves on water to determine surface
tension and viscosity of the given water sample.

\textbf{List of apparatus}

\begin{center}\includegraphics[width=0.53\linewidth]{figures/IPhO_2015_Q5_fig1.jpg}\end{center}

\begin{center}
\begin{tabular}{|c|p{0.62\linewidth}|}
\hline
{[1]} & Light meter (connected to light sensor assembly) \\ \hline
{[2]} & Light sensor assembly mounted on vernier caliper placed on a screen base \\ \hline
{[3]} & Tablet computer (used as sine wave generator) \\ \hline
{[4]} & Digital multimeter \\ \hline
{[5]} & Vibrator control box \\ \hline
{[6]} & Wooden platform \\ \hline
{[7]} & Track for moving light sensor assembly \\ \hline
{[8]} & DC regulated power supply \\ \hline
{[9]} & Hex key, measuring tape and plastic scale \\ \hline
\end{tabular}
\end{center}

\begin{center}Figure 1: Wooden platform unit\end{center}

% NOTE: in the source this photo is cut off at the bottom of page 1 (label 15
% only half visible); its table cell continues empty on page 2.
\begin{center}\includegraphics[width=0.52\linewidth]{figures/IPhO_2015_Q5_fig2.jpg}\end{center}

\begin{center}
\begin{tabular}{|c|p{0.62\linewidth}|}
\hline
{[10]} & Scale and rider with vibrator position marker \\ \hline
{[11]} & Vibrator assembly \\ \hline
{[12]} & Water tray \\ \hline
{[13]} & Plastic cover \\ \hline
{[14]} & Assembly for adjusting vibrator height \\ \hline
{[15]} & Laser source 2\newline (Wavelength, $\lambda_\mathrm{L} = 635\,\mathrm{nm}$, $1\,\mathrm{nm} = 10^{-9}\,\mathrm{m}$) \\ \hline
{[16]} & Water sample for the experiment \\ \hline
{[17]} & $500\,\mathrm{ml}$ measuring cylinder \\ \hline
\end{tabular}
\end{center}

\begin{center}Figure 2: Vibrator/laser source unit\end{center}

\textbf{Description of apparatus}

\textbf{a)\quad Tablet computer as sine wave generator}

\begin{center}\includegraphics[width=0.3\linewidth]{figures/IPhO_2015_Q5_fig3.png}\end{center}

\begin{center}
\begin{tabular}{|p{0.85\linewidth}|}
\hline
{[18]}: Power Switch \\ \hline
{[19]}: Volume up \\ \hline
{[20]}: Volume down \\ \hline
{[21]}: Charging port \\ \hline
{[22]}: Socket for Audio connector pin of cable coming from vibrator control box[5] \\ \hline
\end{tabular}
\end{center}

\begin{center}Figure 3: Switches and sockets of the tablet\end{center}

\begin{tabular}{@{}lp{0.8\linewidth}@{}}
Note & \textbullet\ Keep the tablet always charging. \\
 & \textbullet\ Gently press the power switch \textit{once} to display initial screen. \\
 & \textbullet\ Keep the output volume at maximum using the “Volume up” button[19]. \\
\end{tabular}

\begin{center}
\begin{tabular}{|m{0.22\linewidth}|m{0.16\linewidth}|m{0.22\linewidth}|m{0.2\linewidth}|}
\hline
\includegraphics[width=\linewidth]{figures/IPhO_2015_Q5_fig4.jpg} &
\centering\includegraphics[width=0.35\linewidth]{figures/IPhO_2015_Q5_fig5.png}\par\raggedright Touch and slide the icon[23] to unlock &
\includegraphics[width=\linewidth]{figures/IPhO_2015_Q5_fig6.jpg} &
\centering\includegraphics[width=0.3\linewidth]{figures/IPhO_2015_Q5_fig7.png}\par\raggedright Tap the icon [24] to start the sine wave generator \tabularnewline \hline
\end{tabular}
\end{center}

\begin{center}Figure 4: Initial screens of the tablet\end{center}

\begin{center}\includegraphics[width=0.42\linewidth]{figures/IPhO_2015_Q5_fig8.jpg}\end{center}

\begin{center}
\begin{tabular}{|p{0.6\linewidth}|}
\hline
{[25]}: Waveform selector (always keep at “SIN”) \\ \hline
{[26]}: Amplitude slider \\ \hline
{[27]}: Frequency slider \\ \hline
{[28]}: Frequency value field [Hz] \\ \hline
{[29]}: Application status indicator/switch\newline
“OFF” - the sine wave generator is OFF\newline
“ON” - the sine wave generator is ON \\ \hline
\end{tabular}
\end{center}

\begin{center}Figure 5: The sine wave generator application\end{center}

\begin{center}\includegraphics[width=0.34\linewidth]{figures/IPhO_2015_Q5_fig9.jpg}\end{center}

\textit{To vary frequency}\\
\textbullet\ Tap frequency value field [28] (Fig. 5) to reveal number pad\\
\textbullet\ Use backspace button [30] to erase frequency value\\
\textbullet\ Enter the required frequency, and press “Finished” button[31]

\begin{center}Figure 6: Screen showing number pad to enter frequency value\end{center}

\textit{To vary amplitude}\\
\textbullet\ Use amplitude slider [26] on tablet screen or the variable knob [33]
on vibrator control box [5] to vary output amplitude.

\textbf{b) Vibrator control box, digital multimeter, DC regulated power supply
and their connections}

\begin{center}\includegraphics[width=0.21\linewidth]{figures/IPhO_2015_Q5_fig10.jpg}\end{center}

\begin{center}
\begin{tabular}{|p{0.6\linewidth}|}
\hline
{[32]}: Sockets to connect cables from multimeter \\ \hline
{[33]}: Knob for varying amplitude of the sine wave \\ \hline
{[34]}: Socket for pin of the cable from the vibrator assembly \\ \hline
{[35]}: USB pin to be connected to DC regulated power supply \\ \hline
{[36]}: Audio pin to be connected to the tablet \\ \hline
\end{tabular}
\end{center}

\begin{center}Figure 7: Vibrator control box[5]\end{center}

\begin{center}\includegraphics[width=0.29\linewidth]{figures/IPhO_2015_Q5_fig11.jpg}\end{center}

\begin{center}
\begin{tabular}{|p{0.6\linewidth}|}
\hline
{[37]}: Vibrator strip \\ \hline
{[38]}: Pin of the cable from vibrator assembly \\ \hline
\end{tabular}
\end{center}

\begin{center}Figure 8: Vibrator assembly[11]\end{center}

\begin{center}\includegraphics[width=0.12\linewidth]{figures/IPhO_2015_Q5_fig12.jpg}\end{center}

\begin{center}
\begin{tabular}{|p{0.6\linewidth}|}
\hline
{[39]}: AC/DC selector switch \\ \hline
{[40]}: Range selector knob \\ \hline
{[41]}: Input sockets \\ \hline
\end{tabular}
\end{center}

\begin{center}Figure 9: Digital multimeter[4]\end{center}

\begin{center}\includegraphics[width=0.3\linewidth]{figures/IPhO_2015_Q5_fig13.jpg}\end{center}

\begin{center}Figure 10: Laser source 2 [15] (mounted on a metal block) with connector [42]\end{center}

\begin{center}\includegraphics[width=0.23\linewidth]{figures/IPhO_2015_Q5_fig14.jpg}\end{center}

\begin{center}
\begin{tabular}{|p{0.6\linewidth}|}
\hline
{[43]}: Intensity switch (keep on “High” position) \\ \hline
{[44]}: USB socket for USB pin from vibrator control box \\ \hline
{[45]}: Socket for connector from laser source 2 \\ \hline
\end{tabular}
\end{center}

\begin{center}Figure 11: DC regulated power supply[8]\end{center}

\begin{center}\includegraphics[width=0.96\linewidth]{figures/IPhO_2015_Q5_fig15.jpg}\end{center}

\begin{center}
\begin{tabular}{|p{0.24\linewidth}|p{0.16\linewidth}|p{0.18\linewidth}|p{0.26\linewidth}|}
\hline
\centering [36]$\rightarrow$[22] & \centering [38]$\rightarrow$[34] & \centering [41]$\leftrightarrow$[32] & \centering [35]$\rightarrow$[44] and [42]$\rightarrow$[45] \tabularnewline \hline
\end{tabular}
\end{center}

\begin{center}Figure 12: Connections between tablet, vibrator control box and DC regulated power supply\end{center}

\textbf{c) Light sensor assembly and light meter}

\begin{center}\includegraphics[width=0.45\linewidth]{figures/IPhO_2015_Q5_fig16.jpg}\end{center}

\begin{center}
\begin{tabular}{|p{0.6\linewidth}|}
\hline
{[46]}: Circular aperture on the light sensor\newline
{[47]}: Power switch of the light meter\newline
{[48]}: A, B, C – Sensitivity ranges of the light meter \\ \hline
\end{tabular}
\end{center}

\begin{center}Figure 13: Light sensor assembly and light meter\end{center}

\begin{center}
\begin{tabular}{|p{0.3\linewidth}|p{0.3\linewidth}|}
\hline
\includegraphics[width=\linewidth]{figures/IPhO_2015_Q5_fig17.jpg} &
\includegraphics[width=\linewidth]{figures/IPhO_2015_Q5_fig18.jpg} \\ \hline
One jaw of the vernier caliper fits into a slot at the back of the light sensor &
Tighten the screw using the hex key \\ \hline
\end{tabular}
\end{center}

\begin{center}Figure 14: Attaching light sensor assembly\end{center}

\textbf{Initial Adjustments}

\begin{center}
\begin{tabular}{|p{0.16\linewidth}|p{0.18\linewidth}|p{0.55\linewidth}|}
\hline
\includegraphics[width=\linewidth]{figures/IPhO_2015_Q5_fig19.jpg} &
\includegraphics[width=\linewidth]{figures/IPhO_2015_Q5_fig20.jpg} &
\includegraphics[width=\linewidth]{figures/IPhO_2015_Q5_fig21.png} \\ \hline
\centering Figure 15: Removing the right reflector &
Figure 16: Base screws touching the wooden strip &
\centering Figure 17: Correct position of the vibrator strip and black knob for height adjustment \tabularnewline \hline
\end{tabular}
\end{center}

1. Disconnect the laser 1 connector and insert the laser 2 connector into the
socket of the DC regulated power supply. Note: Laser 2 has been already
adjusted for a specific angle of incidence. Do not touch the laser source!

2. Remove the right reflector used in E-I by turning the bolt under the wooden
platform (Fig. 15).

3. Remove the screen used in E-I and insert the light sensor assembly into the
screen base. Place the screen base between the guiding rails of the track.

4. Position the wooden platform [6] with its base screws touching the wooden
strip attached to the working table (Fig. 16).

5. Raise the side flap of the plastic cover on the vibrator/laser source unit.
Pour exactly $500\,\mathrm{ml}$ of the water sample into the tray [12] using the
measuring cylinder [17].

6. Switch on the laser. Locate the reflected laser spot on the light sensor. As
you move the light sensor assembly back and forth along the track, the laser
spot should move vertically and not at an angle to the vertical. Minor lateral
adjustment of the wooden platform and vertical movement of light sensor
assembly will allow you to get the laser spot exactly on the aperture. The
intensity shown by the light meter will be maximum, if the centre of the laser
spot coincides with the centre of the aperture,.

7. The vibrator strip has already been arranged in the correct vertical
position. Do NOT change the black knob of the height adjustment assembly [14]
(Fig. 17).

8. The vibrator assembly can be moved back and forth horizontally. Vibrator
position marker indicates the position of the assembly on the scale [10].

9. While recording data, keep the flap of the plastic cover lowered in order to
protect the water surface from air currents.

{\large\textbf{Experiment}}

\textbf{Part C: Measurement of angle $\boldsymbol{\theta}$ between the laser beam
and the water surface}

\begin{center}\includegraphics[width=0.78\linewidth]{figures/IPhO_2015_Q5_fig22.png}\end{center}

\begin{center}Figure 18: Measurement of angle $\theta$\end{center}

\begin{center}
\begin{tabular}{|c|p{0.74\linewidth}|c|}
\hline
Tasks & \centering Description & Marks \\ \hline
C1 & Move the light sensor assembly in suitable steps along the track. Note down
the X–displacement of the assembly and the corresponding Y-displacement of the
laser spot. Record your readings in Table C1. (Select appropriate range in the
light meter.) & \textbf{1.0} \\ \hline
C2 & Plot a suitable graph (label it Graph C1) and determine the angle $\theta$
in degrees from its slope. & \textbf{0.6} \\ \hline
\end{tabular}
\end{center}

\textbf{Part D: Determination of surface tension $\boldsymbol{\sigma}$ of the
given water sample}

From diffraction theory it can be shown that
\begin{equation}
k = \tfrac{2\pi}{\lambda_L} sin\theta\ sin\gamma \tag{1}
\end{equation}
where, $k = {}^{2\pi}/_{\lambda_w}$ is the wave number of the surface tension
waves,

\hspace*{2em}$\lambda_\mathrm{w}$ and $\lambda_\mathrm{L}$ being the
wavelengths of the surface tension waves and the laser respectively.

\hspace*{2em}The angle $\gamma$ is the angular distance between the central
maximum and the first-order maximum (Fig. 19).

The vibration frequency (\textit{f}) of the waves is related to the wave number
\textit{k} by
\begin{equation}
\omega = \sqrt{\tfrac{\sigma}{\rho} k^q} \tag{2}
\end{equation}
where, $\omega = 2\pi f$, $\rho$ is the density of the water and $q$ is an
integer.

\begin{center}\includegraphics[width=0.91\linewidth]{figures/IPhO_2015_Q5_fig23.png}\end{center}

\begin{center}Figure 19: Schematic diagram of the apparatus\end{center}

1.\quad Fix the light sensor assembly (using the tightening bolt at the screen
base) at the position shown in Fig. 1. Select the appropriate range on the
light meter.

\begin{center}
\begin{tabular}{|c|p{0.74\linewidth}|c|}
\hline
Task & \centering Description & Marks \\ \hline
D1 & Measure the length $l_1$ between the light sensor aperture and outer edge
of the water tray. You will see a line where the laser strikes the water
surface. The centre of this line is the point of incidence of the laser.
Measure $l_2$, the distance of this point from the edge. Obtain $L$. Record it
on your answersheet. & \textbf{0.3} \\ \hline
\end{tabular}
\end{center}

2.\quad Set the vibrator position marker at $7.0\,\mathrm{cm}$ mark on the
horizontal scale [10].

3.\quad Set the sine wave frequency to $60\,\mathrm{Hz}$ and adjust its
amplitude such that the first- and second-order maxima of the diffraction
pattern are clearly visible (Fig. 19 inset).

\begin{center}
\begin{tabular}{|c|p{0.74\linewidth}|c|}
\hline
Tasks & \centering Description & Marks \\ \hline
D2 & Measure the distance between the second-order maximum above and below the
central maximum. Hence calculate $x_1$. Record your observations in Table D1.
Repeat this by increasing the frequencies in appropriate steps. & \textbf{2.8} \\ \hline
D3 & Identify the appropriate variables for a suitable graph whose slope would
give the value of $q$. Enter the variable values in Table D2. Plot the graph to
find $q$ (label it Graph D1). Write down equation 2 with the appropriate
integer value of $q$. & \textbf{0.9} \\ \hline
D4 & From the equation 2, identify the appropriate variables for a suitable
graph whose slope would give the value of $\sigma$. Enter the variable values
in Table D3. Plot the graph to determine $\sigma$ (label it Graph D2).
($\rho = 1000\,\mathrm{kg.m^{-3}}$). & \textbf{1.2} \\ \hline
\end{tabular}
\end{center}

\textbf{Part E: Determination of the attenuation constant,
$\boldsymbol{\delta}$ and the viscosity of the liquid, $\boldsymbol{\eta}$}

The surface tension waves are damped due to the viscosity of water. The wave
amplitude, $h$, decreases exponentially with the distance, $s$, measured from
the vibrator,
\begin{equation}
h = h_0 e^{-\delta s} \tag{3}
\end{equation}
where, $h_0$ is the amplitude at the vibrator position and $\delta$ is the
attenuation constant.

Experimentally, amplitude $h_0$ can be related to the voltage
($V_\mathrm{rms}$) applied to the vibrator assembly as,
\begin{equation}
h_0 \propto (V_{rms})^{0.4} \tag{4}
\end{equation}
The attenuation constant is related to the viscosity of the liquid as
\begin{equation}
\delta = \tfrac{8}{3}\tfrac{\pi\eta f}{\sigma} \tag{5}
\end{equation}
where, $\eta$ is the viscosity of the liquid.

\hspace*{1em}1.\quad Set the vibrator position marker at $8.0\,\mathrm{cm}$.

\hspace*{1em}2.\quad Adjust the frequency to $100\,\mathrm{Hz}$.

\hspace*{1em}3.\quad Adjust the light sensor using the vernier caliper such
that the first-order maximum of the diffraction pattern falls on the aperture.

\hspace*{1em}4.\quad Adjust the amplitude of sine wave ($V_\mathrm{rms}$) such
that the reading in the light meter is 100 on range A. Note down
$V_\mathrm{rms}$ corresponding to the light meter reading.

\hspace*{1em}5.\quad Move the vibrator away from the point of incidence of the
laser in steps of $0.5\,\mathrm{cm}$ and adjust $V_\mathrm{rms}$ to get the
light meter reading 100. Note down corresponding $V_\mathrm{rms}$.

\begin{center}
\begin{tabular}{|c|p{0.74\linewidth}|c|}
\hline
Tasks & \centering Description & Marks \\ \hline
E1 & Record your data for every step in Table E1. & \textbf{1.9} \\ \hline
E2 & Plot a suitable graph (label it Graph E1) and determine the attenuation
constant $\delta$ from its slope. & \textbf{1.0} \\ \hline
E3 & Calculate the viscosity $\eta$ of the given water sample. & \textbf{0.3} \\ \hline
\end{tabular}
\end{center}
